Une grille rectangulaire comporte $m$ lignes et $n$ colonnes. Toutes les
cases sont initialement blanches.

Une opération consiste à choisir deux lignes et deux colonnes, puis à changer
la couleur des quatre cases situées à leurs intersections.
\begin{enumerate}
  \item Montrer que, dans chaque ligne, la parité du nombre de cases noires
  est invariante.
  \item Montrer de même que, dans chaque colonne, cette parité est invariante.
  \item En déduire que toute configuration accessible possède un nombre pair
  de cases noires dans chaque ligne et dans chaque colonne.
  \item Réciproquement, considérer une configuration possédant cette
  propriété. Pour chaque case noire $(i, j)$, avec $i < m$ et $j < n$,
  effectuer l'opération sur les lignes $i, m$ et les colonnes $j, n$.
  \item Montrer que cette procédure permet d'effacer toutes les cases noires.
  \item Conclure qu'une configuration est accessible si et seulement si chaque
  ligne et chaque colonne contient un nombre pair de cases noires.
\end{enumerate}
